\Needspace{12\baselineskip}
\begin{center}
\textit{Formen mit quadratischer Determinante.}
\end{center}

\Needspace{8\baselineskip}
\Article{206}
\textsc{Aufgabe.} \textit{Für eine gegebene Form $(a,b,c)$ mit quadratischer Determinante $hh$, wobei $h$ deren positive Wurzel bezeichnet, eine eigentlich äquivalente Form $(A,B,C)$ finden, in der $A$ zwischen $0$ und $2h-1$ einschließlich liegt und $B$ $=h$ sowie $C=0$ gilt.}

\textit{Lösung.} I. Wegen $hh=bb-ac$ ist $(h-b):a=c:-(h+b)$. Diesem Verhältnis sei $\beta:\delta$ gleich, wobei $\beta$ und $\delta$ teilerfremd sind. Wir bestimmen $\alpha$, $\gamma$ so, dass $\alpha\delta-\beta\gamma=1$ gilt, was möglich ist. Durch die Substitution $\alpha,\beta,\gamma,\delta$ gehe die Form $(a,b,c)$ in $(a',b',c')$ über, die somit zu ihr eigentlich äquivalent ist. Wir erhalten
\[
\begin{aligned}
b'&=a\alpha\beta+b(\alpha\delta+\beta\gamma)+c\gamma\delta\\
&=(h-b)\alpha\delta+b(\alpha\delta+\beta\gamma)-(h+b)\beta\gamma\\
&=h(\alpha\delta-\beta\gamma)=h\\
c'&=a\beta\beta+2b\beta\delta+c\delta\delta\\
&=(h-b)\beta\delta+2b\beta\delta-(h+b)\beta\delta=0
\end{aligned}
\]
Liegt also außerdem $a'$ bereits zwischen $0$ und $2h-1$, so erfüllt die Form $(a',b',c')$ alle Bedingungen.

II. Liegt dagegen $a'$ außerhalb der Grenzen $0$ und $2h-1$, so sei $A$ der kleinste positive Rest von $a'$ modulo $2h$, der offenbar zwischen diesen Grenzen liegt; wir setzen $A-a'=2hk$. Dann geht die Form $(a',b',c')$, \textit{d.h.} $(a',h,0)$, durch die Substitution $1,0,k,1$ in $(A,h,0)$ über, die zu $(a',b',c')$, $(a,b,c)$ eigentlich äquivalent ist und alle Bedingungen erfüllt. --- Offenbar geht die Form $(a,b,c)$ in $(A,h,0)$ über, und zwar durch die Substitution $\alpha+\beta k,\beta,\gamma+\delta k,\delta$.

\textit{Beispiel.} Gegeben sei die Form $(27,15,8)$ mit der Determinante $=9$. Hier ist $h=3$; den Verhältnissen $-12:27=8:-18$ ist in kleinsten Zahlen das Verhältnis $4:-9$ gleich. Setzen wir also $\beta=4$, $\delta=-9$, $\alpha=-1$, $\gamma=2$, so wird die Form $(a',b',c')$ zu $(-1,3,0)$, die in $(5,3,0)$ übergeht, und zwar durch die Substitution $1,0,1,1$. Das ist somit
\SourceLeaf{206}
die gesuchte Form, und die gegebene geht durch die eigentliche Substitution $3,4,-7,-9$ in sie über.

Solche Formen $(A,B,C)$, bei denen $C=0$, $B=h$ ist und $A$ zwischen $0$ und $2h-1$ liegt, nennen wir \textit{reduzierte Formen}. Sie müssen also von den reduzierten Formen mit negativer oder positiver nichtquadratischer Determinante sorgfältig unterschieden werden.

\Needspace{30\baselineskip}
\Article{207}
\textsc{Satz.} \textit{Zwei reduzierte Formen $(a,h,0)$, $(a',h,0)$, die nicht identisch sind, können nicht eigentlich äquivalent sein.}

\textit{Beweis.} Nehmen wir an, sie seien eigentlich äquivalent, und die erste gehe durch die eigentliche Substitution $\alpha,\beta,\gamma,\delta$ in die zweite über. Dann haben wir vier Gleichungen:
\begin{align*}
a\alpha\alpha+2h\alpha\gamma&=a'\tag*{. . . . . . . . . [1]}\\
a\alpha\beta+h(\alpha\delta+\beta\gamma)&=h\tag*{. . . . . . . [2]}\\
a\beta\beta+2h\beta\delta&=0\tag*{. . . . . . . . . . [3]}\\
\alpha\delta-\beta\gamma&=1\tag*{. . . . . . . . . . . [4]}
\end{align*}
Multiplizieren wir die zweite Gleichung mit $\beta$, die dritte mit $\alpha$ und subtrahieren, so folgt $-h(\alpha\delta-\beta\gamma)\beta=\beta h$ oder wegen [4] $-\beta h=\beta h$, also notwendig $\beta=0$. Daher ergibt [4] $\alpha\delta=1$ und $\alpha=\pm1$. Aus [1] folgt dann $a\pm2\gamma h=a'$. Diese Gleichung kann nur gelten, wenn $\gamma=0$ ist (da sowohl $a$ als auch $a'$ nach Voraussetzung zwischen $0$ und $2h-1$ liegen), \textit{d.h.}, wenn $a=a'$ ist oder die Formen $(a,h,0)$, $(a',h,0)$ identisch sind, entgegen der Voraussetzung.

Damit lassen sich die folgenden Aufgaben, die bei nichtquadratischen Determinanten viel größere Schwierigkeiten bereiteten, leicht lösen.

I. \textit{Für zwei gegebene Formen $F$, $F'$ mit derselben quadratischen Determinante untersuchen, ob sie eigentlich äquivalent sind.} Wir suchen zwei reduzierte Formen, die zu $F$, $F'$ der Reihe nach eigentlich äquivalent sind. Sind diese identisch, so sind die gegebenen Formen eigentlich äquivalent; andernfalls nicht.

II. \textit{Unter denselben Voraussetzungen untersuchen, ob sie uneigentlich äquivalent sind.} Die einer der gegebenen Formen, \textit{z.B.} $F$, entgegengesetzte Form sei $G$. Ist diese zu $F'$ eigentlich äquivalent, so sind $F$ und $F'$ uneigentlich äquivalent, und umgekehrt.
